% MA 214 Discussion 3 -- covers Lectures 4 and 5.
%
% Student handout. Page 1 is the concept review; the question set runs from page 2.
% Compile from inside this folder:  pdflatex Discussion3.tex

\documentclass[11pt]{article}
\usepackage[margin=1in]{geometry}
\usepackage{parskip}
\usepackage{fancyhdr}
\usepackage{array}
\usepackage{booktabs}
\usepackage{enumitem}
\usepackage{amsmath}
\usepackage{tabularx}
\usepackage{listings}

\pagestyle{fancy}
\fancyhf{}
\cfoot{\thepage}
\rhead{Discussion 3}
\lhead{MA 214: Applied Statistics}
\renewcommand{\headrulewidth}{.4mm}

\newcommand{\blankline}[1]{\underline{\hspace{#1}}}
\newcolumntype{L}{>{\raggedright\arraybackslash}X}

\lstset{
  basicstyle=\ttfamily\small,
  columns=fullflexible,
  frame=single,
  framerule=0.4pt,
  breaklines=true,
  showstringspaces=false,
  aboveskip=8pt,
  belowskip=8pt
}

% Section header: bold title over a hairline rule.
\newcommand{\parttitle}[1]{%
  \par\medskip\noindent
  {\large\bfseries #1}\par
  \vspace{-.5em}\noindent\rule{\linewidth}{0.4pt}\par\smallskip}

% Writing space.
\newcommand{\answerspace}[1]{\vspace{#1}}

% Flush-left question list: the label runs in at the margin and the body wraps
% back to the margin, so nothing is pushed far to the right.
\newlist{questions}{enumerate}{1}
\setlist[questions]{label=\textbf{Question \arabic*.}, wide=0pt, itemsep=1.4em, topsep=.8em}

\newlist{parts}{enumerate}{1}
\setlist[parts]{label=\textbf{(\Alph*)}, leftmargin=1.6em, labelsep=.45em,
  itemsep=.7em, topsep=.5em}

\begin{document}

\noindent\textbf{Name:}\blankline{2.3in}\hfill\textbf{BUID:}\blankline{1.6in}

\begin{center}
  {\Large\bfseries Discussion 3: From One Predictor to Several}
\end{center}

\noindent\textbf{Goals:} \textbf{(1)} say what $r$ and $R^2$ do and do not measure,
\textbf{(2)} interpret a coefficient as a claim made \emph{holding the other predictors fixed},
\textbf{(3)} explain why a slope changes when a second predictor is added, \textbf{(4)} read an
indicator-variable model and find its reference level, \textbf{(5)} choose between models with
adjusted $R^2$ rather than $R^2$.

\parttitle{Part 1: Concept Review}

{\small

\noindent Everything in Part 2 can be answered from this page. Keep it in front of you.

\begin{enumerate}[label=\textbf{\arabic*.}, wide=0pt, itemsep=.55em, topsep=.5em]

\item \textbf{Correlation and $R^2$.} The correlation $r$ measures the direction and strength of a
\emph{linear} relationship, and $-1 \le r \le 1$. For simple linear regression $R^2 = r^2$, read as
``the fraction of the variability in $y$ explained by the model.'' Neither one measures steepness:
$\hat\beta_1 = r \cdot s_y / s_x$, so two data sets can share an $R^2$ and have completely
different slopes. Neither tells you the model is \emph{correct}, only how tightly the points hug
whatever line was fit.

\item \textbf{Prediction, residuals, extrapolation.} A fitted model gives
$\hat{y} = \hat\beta_0 + \hat\beta_1 x_1 + \cdots + \hat\beta_p x_p$, and the residual is
$e_i = y_i - \hat{y}_i$. Predicting inside the observed range of the predictors is
\textbf{interpolation}; predicting outside it is \textbf{extrapolation}, and the arithmetic gives
no warning when you have left the data behind.

\item \textbf{Multiple regression: every coefficient is a conditional claim.} In
$\hat{y} = \hat\beta_0 + \hat\beta_1 x_1 + \hat\beta_2 x_2$, the coefficient $\hat\beta_1$ is the
predicted change in $y$ for a one-unit increase in $x_1$ \textbf{holding $x_2$ fixed}. Drop $x_2$
from the model and $\hat\beta_1$ generally changes, because the simple slope absorbs whatever
$x_2$ was contributing through its association with $x_1$. Neither slope is ``wrong''; they answer
different questions.

\item \textbf{Indicator variables.} A categorical predictor with $k$ levels enters as $k-1$
indicator (0/1) variables. The omitted level is the \textbf{reference level}: the intercept
describes it, and each indicator coefficient is the predicted difference between its level and the
reference, holding the other predictors fixed. Changing the reference level changes the
coefficients but not the fitted values.

\item \textbf{$R^2$ versus adjusted $R^2$.} With $p$ predictors and $n$ observations,
\[
R^2 = 1 - \frac{\text{SSE}}{\text{SST}},
\qquad
R^2_{\text{adj}} = 1 - \frac{\text{SSE}/(n-p-1)}{\text{SST}/(n-1)}
= 1 - (1-R^2)\,\frac{n-1}{n-p-1}.
\]
Adding any predictor, useful or not, can only \emph{raise} $R^2$, so $R^2$ cannot be used to
compare models of different sizes. Adjusted $R^2$ charges a price for each predictor and can go
down, which is what makes it usable for that comparison.

\item \textbf{Four traps.} Reading $R^2$ as a measure of slope, or as evidence the model is right.
Interpreting a multiple-regression coefficient without the phrase ``holding the others fixed.''
Forgetting which level is the reference, so an indicator coefficient is read as a level mean
rather than a difference. Comparing models by $R^2$ when they have different numbers of
predictors.

\end{enumerate}

}

\newpage

\parttitle{Part 2: Question Set}

\noindent\textbf{Scenario.} The bike-share system has 12 stations. For each one you have
\texttt{docks} (number of docks, in tens), \texttt{type} (Campus, Downtown, or Residential), and
\texttt{trips} (daily trips started, in tens). Docks range from 4 to 12.

\begin{center}
\renewcommand{\arraystretch}{1.15}
\begin{tabular}{lcccccccccccc}
\toprule
\textbf{Station} & A & B & C & D & E & F & G & H & I & J & K & L \\
\midrule
\texttt{docks} & 5 & 6 & 7 & 8 & 9 & 10 & 11 & 12 & 4 & 5 & 6 & 7 \\
\texttt{trips} & 12.4 & 12.9 & 14.6 & 15.5 & 25.3 & 25.6 & 27.3 & 28.7 & 7.6 & 9.4 & 9.9 & 11.5 \\
\texttt{type} & C & C & C & C & D & D & D & D & R & R & R & R \\
\bottomrule
\end{tabular}
\end{center}

\noindent (C = Campus, D = Downtown, R = Residential.)

\begin{questions}

%%%%%%%%%%%%%%%%%%%%%%%%%%%%%%%%%%%%%%
\item \textbf{The one-predictor fit.} Regressing \texttt{trips} on \texttt{docks} alone gives

\begin{center}
\renewcommand{\arraystretch}{1.2}
\begin{tabular}{lcc}
\toprule
 & \textbf{Estimate} & \textbf{Std. Error} \\
\midrule
(Intercept) & $-5.041$ & 2.286 \\
\texttt{docks} & $2.902$ & 0.290 \\
\bottomrule
\end{tabular}
\qquad
$r = 0.954$, \quad $R^2 = 0.909$
\end{center}

\begin{parts}
\item Interpret the slope $2.902$ in context, \emph{with units}. Remember what the tens mean.
\answerspace{4em}

\item Interpret $R^2 = 0.909$ in one sentence. Then say what it does \emph{not} tell you.
\answerspace{4em}

\item What does the intercept $-5.041$ predict, and why should you refuse to report it?
\answerspace{3.5em}

\item Two other cities fit the same model to their own stations. Both report $R^2 = 0.909$, but
one has a slope of $1.2$ and the other a slope of $5.8$. Is that possible? Explain using the
relationship between $\hat\beta_1$, $r$, and the two standard deviations.
\answerspace{4.5em}
\end{parts}

%%%%%%%%%%%%%%%%%%%%%%%%%%%%%%%%%%%%%%
\item \textbf{Adding a second predictor.} Now add \texttt{type} to the model:

\begin{center}
\renewcommand{\arraystretch}{1.2}
\begin{tabular}{lcc}
\toprule
 & \textbf{Estimate} & \textbf{Std. Error} \\
\midrule
(Intercept) & $6.245$ & 0.627 \\
\texttt{docks} & $1.170$ & 0.092 \\
\texttt{type}Downtown & $8.195$ & 0.448 \\
\texttt{type}Residential & $-3.080$ & 0.269 \\
\bottomrule
\end{tabular}
\qquad
$R^2 = 0.998$
\end{center}

\begin{parts}
\item The \texttt{docks} slope fell from $2.902$ to $1.170$. Look back at the data table and
explain, in terms of \texttt{type}, why the one-predictor slope was so much larger.
\answerspace{5em}

\item Write the interpretation of $1.170$ as a full sentence, including the conditional phrase.
\answerspace{4em}

\item The city can add 10 docks to one existing station. Which of the two slopes answers ``how
many more daily trips should we expect?'' Which one answers ``if I pick two stations at random and
one has 10 more docks, how much busier is it?'' Say which is which.
\answerspace{4.5em}

\item Is the one-predictor slope \emph{wrong}? Answer in one sentence.
\answerspace{3em}
\end{parts}

%%%%%%%%%%%%%%%%%%%%%%%%%%%%%%%%%%%%%%
\item \textbf{Reading the indicator variables.} Use the model from Question 2.

\begin{parts}
\item \texttt{type} has three levels but only two coefficients appear. How many indicators does R
create, and which level is the reference?
\answerspace{3em}

\item Fill in the predicted daily trips at \texttt{docks} $= 8$ for each type.

\begin{center}
\renewcommand{\arraystretch}{1.7}
\begin{tabular}{lcc}
\toprule
\textbf{Type} & \textbf{Arithmetic} & \textbf{Predicted \texttt{trips}} \\
\midrule
Campus & \blankline{2.1in} & \blankline{.9in} \\
Downtown & \blankline{2.1in} & \blankline{.9in} \\
Residential & \blankline{2.1in} & \blankline{.9in} \\
\bottomrule
\end{tabular}
\end{center}

\item Interpret $8.195$ in context. Your sentence must name what it is being compared to.
\answerspace{3.5em}

\item Suppose you refit with Downtown as the reference level. Which of these change and which do
not: the two indicator coefficients, the intercept, the \texttt{docks} slope, $R^2$, the predicted
value for a Campus station with 8 tens of docks?
\answerspace{4.5em}
\end{parts}

%%%%%%%%%%%%%%%%%%%%%%%%%%%%%%%%%%%%%%
\item \textbf{Choosing a model.} The city repeats the study with $n = 60$ stations and fits three
nested models.

\begin{center}
\renewcommand{\arraystretch}{1.5}
\begin{tabular}{lccc}
\toprule
\textbf{Model} & \textbf{Predictors $p$} & \textbf{$R^2$} & \textbf{$R^2_{\text{adj}}$} \\
\midrule
1: \texttt{docks} & 1 & 0.740 & 0.7355 \\
2: $+$ \texttt{type} & 3 & 0.810 & 0.7998 \\
3: $+$ 3 more predictors & 6 & 0.815 & \blankline{1in} \\
\bottomrule
\end{tabular}
\end{center}

\begin{parts}
\item Compute the missing $R^2_{\text{adj}}$ for Model 3, using the formula in the review.
\answerspace{4em}

\item Which model would you report, and why?
\answerspace{3.5em}

\item Going from Model 2 to Model 3, $R^2$ went up but $R^2_{\text{adj}}$ went down. Explain how
both can happen at once.
\answerspace{4em}

\item A classmate says: ``Model 3 has the highest $R^2$, so it fits best, so it will predict new
stations best.'' Two separate things are wrong with that sentence. Name both.
\answerspace{4.5em}
\end{parts}

%%%%%%%%%%%%%%%%%%%%%%%%%%%%%%%%%%%%%%
\item \textbf{Find the bugs.} A classmate fits the model with \texttt{type} and then predicts. The
code runs with no error message and prints plausible-looking numbers. It is wrong in
\textbf{three} separate ways.

\begin{lstlisting}[language=R]
st$type_num <- as.numeric(st$type)          # Campus = 1, Downtown = 2, Residential = 3

fit1 <- lm(trips ~ docks, data = st)
fit2 <- lm(trips ~ docks + type_num, data = st)

c(summary(fit1)$r.squared, summary(fit2)$r.squared)

predict(fit2, newdata = data.frame(docks = 30, type_num = 2))
\end{lstlisting}

\begin{parts}
\item Bug 1 is \texttt{type\_num}. What does the model now assume about the three types, and why
is that assumption indefensible here? Look at the coefficients in Question 2 before you answer.
\answerspace{5em}

\item Bug 2 is the comparison on the fourth line. Why can those two numbers not settle which model
is better, and what should be compared instead?
\answerspace{4em}

\item Bug 3 is in the \texttt{predict} call. It returns a number without complaint. What is wrong
with that number?
\answerspace{4em}

\item Rewrite the three lines that need fixing.
\answerspace{5em}
\end{parts}

\end{questions}

\end{document}
